\newcommand{\figuraMoleculasSuperficieEsfera}{%
\begin{figure}[htbp]
    \centering
    \begin{tikzpicture}[>=Latex,font=\small]
        % Corte de la esfera y capa molecular superficial.
        \path[fill=orange!18,even odd rule]
            (3.35,3.05) circle (2.62)
            (3.35,3.05) circle (2.15);
        \draw[very thick,blue!70!black] (3.35,3.05) circle (2.62);
        \draw[dashed,orange!80!black,thick] (3.35,3.05) circle (2.15);
        \node[font=\bfseries] at (3.35,6.08)
            {Volumen esférico con $N$ moléculas};

        \foreach \x/\y in {
            2.05/1.38,3.15/0.98,4.25/1.38,1.45/2.25,2.45/2.10,
            3.45/1.78,4.55/2.18,5.10/3.05,1.18/3.25,2.15/3.02,
            3.20/2.85,4.15/3.18,2.00/4.00,3.05/3.82,4.10/4.15,
            2.65/4.92,3.75/4.72}
        {\fill[blue!65] (\x,\y) circle (0.085);}

        \foreach \x/\y in {
            1.45/1.05,2.25/0.58,3.25/0.43,4.30/0.70,5.10/1.35,
            5.72/2.32,5.85/3.35,5.42/4.28,4.62/5.08,3.55/5.55,
            2.45/5.40,1.48/4.92,0.88/4.08,0.73/2.95,0.92/1.92}
        {\fill[red!75] (\x,\y) circle (0.105);}

        \draw[->,very thick,blue!70!black]
            (3.35,3.05)--node[above]{$R$}(5.52,4.25);
        \draw[<->,very thick,red!70!black]
            (5.50,3.05)--node[above,sloped]{$L_{\mathrm{mol}}$}(5.97,3.05);
        \draw[->,thick,red!70!black]
            (0.55,5.38)--(1.35,4.78);
        \node[red!70!black,align=center,anchor=east] at (0.52,5.40)
            {capa superficial\\con $N_S$ moléculas};

        % Ampliación local de la capa superficial.
        \draw[gray!60,dashed,thick] (5.05,4.65)--(7.35,5.52);
        \draw[gray!60,dashed,thick] (5.72,3.82)--(7.35,0.82);
        \draw[very thick,rounded corners=3pt,fill=blue!3]
            (7.35,0.82) rectangle (13.25,5.52);
        \fill[orange!18] (10.30,0.82) rectangle (11.15,5.52);
        \draw[dashed,orange!80!black,thick]
            (10.30,0.82)--(10.30,5.52);
        \draw[very thick,blue!70!black]
            (11.15,0.82)--(11.15,5.52);

        \node[font=\bfseries] at (10.30,5.88)
            {Detalle de la superficie};
        \node[text=black!65] at (8.70,1.12) {interior};
        \node[text=black!65] at (12.20,1.12) {exterior};

        \foreach \x/\y in {
            7.75/1.45,8.50/2.10,9.35/1.55,7.92/3.08,8.82/3.52,
            9.62/2.72,7.72/4.42,8.68/4.75,9.55/4.18}
        {\fill[blue!65] (\x,\y) circle (0.09);}

        \foreach \x/\y in {
            10.48/1.55,10.88/2.30,10.55/3.18,10.92/4.05,10.52/4.82}
        {\fill[red!75] (\x,\y) circle (0.11);}

        \draw[<->,very thick,red!70!black]
            (10.30,5.10)--node[above]{$L_{\mathrm{mol}}$}(11.15,5.10);
        \draw[very thick,red!70!black,-{Latex[length=3mm]}]
            (10.72,3.18)--(11.82,3.18);
        \draw[very thick,red!70!black,-{Latex[length=3mm]}]
            (11.82,2.55)--(10.72,2.55);
        \node[align=center,red!70!black] at (12.18,3.72)
            {moléculas que\\entran y salen};
    \end{tikzpicture}
    \caption{Moléculas asociadas con la superficie de una partícula esférica.
    Solo las moléculas situadas en una capa de espesor
    $L_{\mathrm{mol}}$ pueden entrar o salir de la región superficial. El área
    determina que $N_S\propto N^{2/3}$ y, por estadística de conteo,
    $\Delta N_S\propto N^{1/3}$.}
    \label{fig:taller-moleculas-superficie}
\end{figure}%
}
